\chapter{Chirurgia wzdłuż węzła i niezmienniki czterorozmaitości}
\label{ch:tg-knot-surgery}

W rozdziałach o macierzach Seiferta pokazaliśmy, jak wielomian
Alexandera koduje własności trójwymiarowego uzupełnienia węzła.
Teraz wyjaśnimy, jak ten sam wielomian pozwala rozróżniać
\emph{struktury gładkie} na czterorozmaitościach. Podstawowym
narzędziem jest konstrukcja Fintushela--Sterna, zwana
\emph{chirurgią wzdłuż węzła} (\emph{knot surgery}).

Istotne jest odróżnienie tej operacji od chirurgii Dehna.
W chirurgii Dehna zastępujemy torus pełny w rozmaitości
trójwymiarowej. Tutaj usuwamy otoczenie torusa $T^2$ w
czterorozmaitości, a w jego miejsce wstawiamy iloczyn okręgu
z zewnętrzem węzła. Operacja zachowuje wiele klasycznych
niezmienników rozmaitości, może jednak zmienić jej typ
dyfeomorfizmu.

\section{Dane geometryczne i meridian torusa}
\label{sec:tg-ks-data}

\begin{definicja}[Torus o trywialnej wiązce normalnej]
Niech $X$ będzie gładką, zamkniętą, spójną i zorientowaną
czterorozmaitością, a $T\subset X$ gładko zanurzonym,
zorientowanym torusem. Zakładamy, że jego wiązka normalna
jest trywialna jako zorientowana wiązka rzeczywista
rzędu dwa. Wybieramy trywializację otoczenia
$\nu T\cong T^2\times D^2$. Na torusie brzegowym
$\partial\nu T\cong T^3$ oznaczamy przez $a,b$ dwie
klasy stanowiące dodatnio zorientowaną bazę
$H_1(T;\mathbb Z)$, a przez $m_T$ klasę okręgu
$\{*\}\times\partial D^2$ z orientacją brzegową
dysku normalnego. Orientację dysku wybieramy tak,
aby orientacja produktowa torusa i dysku była
orientacją $X$.
\end{definicja}

Dla zorientowanego torusa w orientowanej czterorozmaitości
warunek trywialności wiązki normalnej jest równoważny
$T\cdot T=0$. Rzeczywiście zorientowane wiązki
rzeczywiste rzędu dwa nad $T^2$ są klasyfikowane
przez klasę Eulera w $H^2(T^2;\mathbb Z)\cong\mathbb Z$,
a jej wartość na klasie fundamentalnej jest
liczbą samoprzecięcia torusa. Ta obserwacja
wyjaśnia geometryczną treść założenia $T^2=0$;
nie mówi jednak nic o grupie fundamentalnej dopełnienia.

\begin{definicja}[Zewnętrze węzła i długość preferowana]
Niech $K\subset S^3$ będzie zorientowanym węzłem.
Wybieramy domknięte otoczenie tubularne $N(K)$ i kładziemy
$E_K:=S^3\setminus\operatorname{int}N(K)$.
Na $\partial E_K$ oznaczamy przez $\mu_K$ klasę
meridianu zorientowanego tak, aby jego liczba
splecenia z $K$ wynosiła $+1$. Przez $\lambda_K$
oznaczamy klasę długości preferowanej: jest to
prymitywny generator jądra odwzorowania
$H_1(\partial E_K;\mathbb Z)\to H_1(E_K;\mathbb Z)$,
którego obraz w $H_1(N(K);\mathbb Z)$ jest
dodatnim generatorem wyznaczonym przez orientację $K$.
W szczególności $(\mu_K,\lambda_K)$ stanowi bazę
pierwszej homologii torusa brzegowego.
\end{definicja}

\begin{lemat}[Homologia zewnętrza]
\label{lem:tg-ks-exterior}
Dla każdego węzła $K\subset S^3$ zachodzi
$H_0(E_K;\mathbb Z)\cong\mathbb Z$,
$H_1(E_K;\mathbb Z)\cong\mathbb Z$ i
$H_i(E_K;\mathbb Z)=0$ dla $i\ge2$.
Meridian $\mu_K$ reprezentuje generator $H_1(E_K;\mathbb Z)$,
natomiast $[\lambda_K]=0$ w tej grupie.
\end{lemat}

\begin{proof}
Zewnętrze jest homotopijnie równoważne $S^3\setminus K$.
Dualność Alexandera daje
$\widetilde H_i(S^3\setminus K;\mathbb Z)
\cong\widetilde H^{2-i}(S^1;\mathbb Z)$.
W szczególności jedyna niezerowa zredukowana grupa
homologii pojawia się w wymiarze jeden i jest wolna
rangi jeden, czyli izomorficzna z $\mathbb Z$.
Izomorfizm z $\mathbb Z$ jest zadany liczbą splecenia
z $K$, dlatego meridian jest generatorem. Z definicji
długość preferowana ma zerowy obraz w $H_1(E_K)$.
\end{proof}

\begin{wniosek}[Homologia części dołączanej]
\label{cor:tg-ks-piece}
Niech $P_K:=S^1\times E_K$, a $s$ niech oznacza
klasę pierwszego czynnika $S^1$. Wówczas
$H_1(P_K;\mathbb Z)\cong\mathbb Z[s]\oplus\mathbb Z[\mu_K]$,
$H_2(P_K;\mathbb Z)\cong\mathbb Z[s\times\mu_K]$
oraz $H_i(P_K;\mathbb Z)=0$ dla $i\ge3$.
Ponadto $\chi(P_K)=0$.
\end{wniosek}

\begin{proof}
Wszystkie grupy homologii z lematu
\ref{lem:tg-ks-exterior} są wolne, więc wzór Künnetha
nie ma składników $\operatorname{Tor}$. Otrzymujemy
$H_1(S^1\times E_K)\cong H_1(S^1)\oplus H_1(E_K)$
oraz $H_2(S^1\times E_K)\cong H_1(S^1)\otimes H_1(E_K)$.
Grupy homologii w stopniach co najmniej trzy są zerowe.
Wreszcie $\chi(S^1\times E_K)=\chi(S^1)\chi(E_K)=0$.
\end{proof}

\section{Definicja chirurgii i sposób sklejenia}
\label{sec:tg-ks-glue}

\begin{definicja}[Chirurgia wzdłuż węzła Fintushela--Sterna]
\label{def:tg-ks-operation}
Niech $X,T,K$ spełniają warunki poprzedniej sekcji.
Niech $A:=X\setminus\operatorname{int}\nu T$.
Wybieramy dyfeomorfizm $\varphi:\partial P_K\to\partial A$
odwracający orientację brzegową i spełniający na
pierwszych grupach homologii równości
$\varphi_*([s])=a$, $\varphi_*([\mu_K])=b$
oraz $\varphi_*([\lambda_K])=m_T$.
Rozmaitość
\[
X_K:=A\cup_\varphi(S^1\times E_K)
\]
nazywamy \emph{wynikiem chirurgii wzdłuż węzła $K$ na torusie $T$}.
Identyfikację przyjmujemy jako część danych standardowej
operacji: jest ona zgodna z wybranym ramowaniem torusa
i z orientacjami na brzegach. W zastosowaniach
formuły Seiberga--Wittena żądamy dodatkowo, aby $T$
był regularnym włóknem otoczenia typu \emph{cusp},
z ramowaniem pochodzącym z tego rozwłóknienia.
Jest to otoczenie włókna ostrzowego: jego model
uchwytowy otrzymuje się z $D^4$ przez dołączenie
uchwytu indeksu dwa wzdłuż prawoskrętnego
trójlistnika z ramowaniem zero. Samo zawieranie
torusa w takim otoczeniu nie jest wystarczające.
Por. definicja przed twierdzeniem 1.1 w pracy
Fintushela--Sterna, \emph{Knots, links, and 4-manifolds},
Invent. Math. 134 (1998), 363--400,
arXiv:dg-ga/9612014v2, s.~2.
Przez $[T_K]$
oznaczamy klasę torusa pochodzącego od
$S^1\times\mu_K$ po niewielkim przesunięciu do wnętrza $P_K$.
\end{definicja}

\begin{uwaga}[Orientacje i zakres definicji]
Równanie $\varphi_*([\lambda_K])=m_T$ to istota
operacji: meridian usuwanego torusa jest zastępowany
przez długość preferowaną węzła.
Sprawdźmy orientacje przy konwencji, że wektor normalny
skierowany na zewnątrz występuje jako pierwszy.
Orientację $\partial N(K)$ wyznacza uporządkowana
para $(\mu_K,\lambda_K)$, gdyż $N(K)$ ma orientację
iloczynu $K\times D^2$. Stąd $\partial E_K$
ma orientację przeciwną. Dla $P_K=S^1\times E_K$
orientacja brzegu jest przeciwna do orientacji
produktowej $S^1\times\partial E_K$ (przeniesienie
wektora normalnego przez pierwszy czynnik
zmienia znak). Zatem $(s,\mu_K,\lambda_K)$
wyznacza orientację $\partial P_K$. Z kolei
$(a,b,m_T)$ wyznacza orientację $\partial\nu T$,
a nie $\partial A$. Dyfeomorfizm torusów
zadany przez te trzy kierunki istnieje
(wybieramy współrzędne torusowe i odwzorowanie liniowe
w odpowiednich bazach) i odwraca orientację
$\partial P_K$ względem $\partial A$. Sam zapis na $H_1$
nie jest stwierdzeniem, że dowolne sklejenie trójtorusów
o takich własnościach daje zawsze tę samą rozmaitość
gładką. W twierdzeniach poniżej używamy standardowego
sklejenia Fintushela--Sterna przy torusie o ustalonym
ramowaniu. W szczególności nie utożsamiamy bez dowodu
wyników wszystkich możliwych sklejeń.
\end{uwaga}

\begin{center}
\begin{tikzpicture}[>=Stealth,line cap=round,line join=round]
  \node[draw=tgBlue,fill=tgBlue!10,rounded corners=4pt,
        text width=4.0cm,align=center,inner sep=7pt] (left) at (-3.25,0)
        {\small $A=X\setminus\operatorname{int}\nu T$\\
         \footnotesize brzeg $T^3$; klasy $a,b,m_T$};
  \node[draw=tgTeal,fill=tgTeal!11,rounded corners=4pt,
        text width=4.0cm,align=center,inner sep=7pt] (right) at (3.25,0)
        {\small $S^1\times E_K$\\
         \footnotesize brzeg $T^3$; klasy $s,\mu_K,\lambda_K$};
  \draw[-{Stealth[length=2.5mm]},tgPurple,line width=1.8pt] (left)--(right);
  \node[fill=white,text=tgNavy,font=\scriptsize,align=center,inner sep=2pt] at (0,.35)
        {$m_T\longleftrightarrow\lambda_K$};
  \node[text=tgCoral,font=\footnotesize] at (0,-1.12)
        {$X_K=A\cup_\varphi(S^1\times E_K)$};
\end{tikzpicture}
\end{center}

\begin{przyklad}[Węzeł trywialny]
Dla węzła trywialnego $U$ mamy $E_U\cong S^1\times D^2$,
przy czym meridian węzła odpowiada kierunkowi $S^1$,
a długość preferowana ogranicza dysk w $E_U$.
Stąd $P_U\cong T^2\times D^2$. Przy standardowym sklejeniu
$X_U\cong X$. Przykład ten ustala, która krzywa
pełni rolę meridianu otoczenia torusa po wymianie.
\end{przyklad}

\section{Homologia, charakterystyka Eulera i sygnatura}
\label{sec:tg-ks-homology}

\begin{stwierdzenie}[Homologia dwóch modeli otoczenia]
\label{prop:tg-ks-boundary-maps}
Przy identyfikacji brzegów z definicji
\ref{def:tg-ks-operation} odwzorowania indukowane
przez włączenia $T^3\hookrightarrow T^2\times D^2$
i $T^3\hookrightarrow S^1\times E_K$ mają izomorficzne
obrazy i jądra we wszystkich stopniach homologii
całkowitej, po utożsamieniu generatorów $a,b,m_T$
z odpowiednio $s,\mu_K,\lambda_K$.
\end{stwierdzenie}

\begin{proof}
W stopniu jeden w pierwszym modelu zachodzi
$a\mapsto a$, $b\mapsto b$, $m_T\mapsto0$.
W drugim modelu mamy $s\mapsto s$,
$\mu_K\mapsto\mu_K$, $\lambda_K\mapsto0$.
Zatem jądro jest generowane w obu przypadkach
przez trzeci element, a dwa pierwsze przechodzą
na bazę grupy $\mathbb Z^2$.

W stopniu dwa homologia torusa brzegowego jest
wolna rzędu trzy, z bazą reprezentowaną przez
$a\times b$, $a\times m_T$, $b\times m_T$.
W pierwszym modelu do generatora $H_2(T^2\times D^2)$
przechodzi wyłącznie $a\times b$.
W drugim modelu analogicznie do generatora $H_2(P_K)$
przechodzi $s\times\mu_K$. Klasa $s\times\lambda_K$
znika z naturalności iloczynu zewnętrznego, ponieważ
$[\lambda_K]=0$ w $H_1(E_K)$. Natomiast klasa
$\mu_K\times\lambda_K$ jest klasą fundamentalną
torusa $\partial E_K$ (z dokładnością do znaku);
jej obraz należy do $H_2(E_K)=0$, więc również znika.
W stopniu zero oba modele są spójne i odwzorowania
są izomorfizmami. W stopniu trzy obie grupy docelowe
są zerowe, więc jądrem jest cała grupa
$H_3(T^3;\mathbb Z)\cong\mathbb Z$.
W stopniach wyższych znikają także grupy źródłowe.
\end{proof}

\begin{wniosek}[Klasyczne niezmienniki operacji]
\label{cor:tg-ks-invariants}
W sytuacji z definicji \ref{def:tg-ks-operation}
rozmaitości $X$ i $X_K$ mają izomorficzne
całkowite grupy homologii, a ponadto
$\chi(X_K)=\chi(X)$ i $\sigma(X_K)=\sigma(X)$.
\end{wniosek}

\begin{proof}
Dwie części dołączane mają takie same grupy homologii,
a stwierdzenie \ref{prop:tg-ks-boundary-maps} identyfikuje
odwzorowania z homologii wspólnego brzegu.
Dla każdego $i$ ciąg Mayera--Vietorisa zawiera krótki ciąg
$0\to\operatorname{coker}j_i\to H_i(X_K;\mathbb Z)
\to\ker j_{i-1}\to0$, gdzie $j_i$ oznacza mapę z
$H_i(T^3;\mathbb Z)$ do $H_i(A;\mathbb Z)\oplus H_i(P_K;\mathbb Z)$.
Odpowiedni ciąg dla $X$ ma izomorficzne jądra i kokernel.
Grupa $\ker j_{i-1}$ jest podgrupą wolnej grupy
$H_{i-1}(T^3;\mathbb Z)$, a więc sama jest wolna;
oba krótkie ciągi rozkładają się jako ciągi grup
abelowych (wybieramy podniesienia elementów bazy
wolnej grupy $\ker j_{i-1}$).
Otrzymujemy stąd izomorfizmy całkowitych grup homologii. Addytywność charakterystyki Eulera
względem rozkładu daje
$\chi(X_K)=\chi(A)+\chi(P_K)-\chi(T^3)$,
$\chi(X)=\chi(A)+\chi(T^2\times D^2)-\chi(T^3)$.
Ponieważ wszystkie trzy składniki zastępowane w tych
równaniach mają charakterystykę Eulera równą zero,
wzór dla $\chi$ jest udowodniony.

Równość sygnatur wynika z addytywności Nowikowa:
wymieniane części mają zerową sygnaturę swoich form
przecięcia, ponieważ ich druga homologia jest
generowana przez torus o zerowym samoprzecięciu.
Dokładniej, generatory pochodzą z homologii brzegu,
więc leżą w jądrze mapy do homologii względnej,
a ta mapa wyznacza formę przecięcia na części
z brzegiem przez dualność Poincarégo--Lefschetza.
Obie formy są zatem zerowe, również na ilorazach
przez radykał.
W tym kroku używamy twierdzenia o addytywności
sygnatur przy sklejaniu wzdłuż wspólnego brzegu;
nie jest ono elementarną konsekwencją samych
liczb Betti'ego.
\end{proof}

\begin{uwaga}[Forma przecięcia a homologia]
Izomorfizm grup $H_2(X_K;\mathbb Z)\cong H_2(X;\mathbb Z)$
nie implikuje automatycznie izometrii form przecięcia.
Sam rachunek $\chi$ i $\sigma$ nie wystarcza do
ustalenia parzystości formy. Poniżej wyprowadzimy
izometrię z odwzorowania zachowującego brzeg.
\end{uwaga}

\begin{proof}[Uzupełnienie: izometria form przecięcia]
Zbudujemy ciągłe odwzorowanie $p:X_K\to X$,
równe identyczności na $A$. Najpierw konstruujemy
$f:E_K\to S^1\times D^2$ takie, że na brzegu
jest standardowym utożsamieniem $\mu_K$ z kierunkiem
$S^1$, a $\lambda_K$ z kierunkiem $\partial D^2$.
Pierwsza składowa istnieje, ponieważ odwzorowania
do $S^1=K(\mathbb Z,1)$ są klasyfikowane przez
$H^1(-;\mathbb Z)$, a klasa generująca $H^1(E_K)$
ogranicza się do klasy przyjmującej wartości $1,0$
na $\mu_K,\lambda_K$. Ograniczenie wybranego
reprezentanta tej klasy jest homotopijne z żądanym
odwzorowaniem brzegowym; własność rozszerzania
homotopii dla pary $(E_K,\partial E_K)$ pozwala
uzyskać dokładną równość na brzegu.
Drugą składową przedłużamy do $D^2$ wprost:
w kołnierzu $\partial E_K\times[0,1]$ mnożymy
zadane odwzorowanie do $\partial D^2$ przez $1-t$,
a poza kołnierzem kładziemy zero.

Otrzymane $f$ jest izomorfizmem na homologii:
wysyła generator $\mu_K$ na generator $H_1(S^1\times D^2)$,
a pozostałe grupy porównujemy z lematem
\ref{lem:tg-ks-exterior}. Iloczyn $\operatorname{id}_{S^1}\times f$
jest izomorfizmem na homologii z naturalności
twierdzenia Künnetha. W standardowych współrzędnych
brzegu jego ograniczenie jest dokładnie przepisanym
sklejeniem do $\partial\nu T$ (utożsamionym z
$\partial A$). Jeżeli wybrany dyfeomorfizm brzegu
różni się od tego modelu, stosujemy homotopię
odwzorowań do $T^3=K(\mathbb Z^3,1)$,
które indukują tę samą mapę na $H_1$,
i własność rozszerzania homotopii dla
$(P_K,\partial P_K)$. Po tej zmianie możemy
skleić z identycznością na $A$.

Naturalność ciągów Mayera--Vietorisa i lemat pięciu
dają, że $p_*$ jest izomorfizmem we wszystkich
stopniach homologii całkowitej. Ponadto $p$ ma
stopień jeden: na małej kuli we wnętrzu $A$
jest identycznością zachowującą orientację, a
obraz $P_K$ leży w $\nu T$ i nie spotyka tej kuli.
Równoważnie, obraz klasy fundamentalnej można
sprawdzić w lokalnej homologii w punkcie tej kuli.
Wobec tego dla $x,y\in H^2(X;\mathbb Z)$
\[
\langle p^*x\smile p^*y,[X_K]\rangle
=\langle x\smile y,p_*[X_K]\rangle
=\langle x\smile y,[X]\rangle.
\]
Z naturalności uniwersalnego twierdzenia
o współczynnikach i lematu pięciu wynika,
że $p^*$ jest izomorfizmem. Aby otrzymać izometrię
w homologii, oznaczmy przez $\operatorname{PD}_M$
izomorfizm $x\mapsto x\frown[M]$. Naturalność
iloczynu z klasą fundamentalną daje
\[
p_*\operatorname{PD}_{X_K}(p^*x)
=\operatorname{PD}_X(x).
\]
Zatem $p_*$ na drugiej homologii jest zgodne
z powyższą izometrią kohomologiczną poprzez
dualność Poincarégo. Po przejściu do ilorazów przez
torsję jest to izometria całkowitych form przecięcia. Ponadto $p_*[T_K]=[T]$, gdyż
generator $s\times\mu_K$ przechodzi na generator
$a\times b$ w $\nu T$.
\end{proof}

\section{Grupa fundamentalna i zachowanie typu homeomorfizmu}
\label{sec:tg-ks-pi1}

\begin{twierdzenie}[Warunek prostej spójności]
\label{thm:tg-ks-pi1}
Niech $X$ będzie gładką, spójną, zamkniętą
czterorozmaitością, a $T\subset X$ torusem o
trywialnej wiązce normalnej i spójnym dopełnieniu.
Jeżeli $\pi_1(X\setminus\operatorname{int}\nu T)=1$,
to dla dowolnego węzła $K\subset S^3$
rozmaitość $X_K$ ma trywialną grupę fundamentalną.
\end{twierdzenie}

\begin{proof}
Oznaczmy $A:=X\setminus\operatorname{int}\nu T$.
Z twierdzenia van Kampena zastosowanego do
$X_K=A\cup_\varphi P_K$ (po poszerzeniu kawałków
przez otwarte kołnierze) dostajemy iloraz wolnego
iloczynu $\pi_1(A)*\pi_1(P_K)$ przez najmniejszą podgrupę normalną zawierającą relacje
$i_A(\varphi_*(\gamma))=i_P(\gamma)$,
gdzie $\gamma\in\pi_1(\partial P_K)$, zaś $i_A,i_P$
są homomorfizmami indukowanymi przez włączenia.
Ponieważ $\pi_1(A)=1$, iloraz ten ma postać
\[
\pi_1(X_K)\cong
\pi_1(P_K)\big/\langle\!\langle
  i_P\bigl(\pi_1(\partial P_K)\bigr)
\rangle\!\rangle.
\]
Obraz grupy brzegu jest generowany przez $s$, $\mu_K$
i $\lambda_K$. Ponieważ
$\pi_1(P_K)\cong\mathbb Z[s]\times\pi_1(E_K)$,
a grupa węzła $\pi_1(E_K)$ jest normalnie generowana
przez meridian $\mu_K$, to relacje $s=1$ i
$\mu_K=1$ już zabijają całą grupę. Otrzymujemy
$\pi_1(X_K)=1$.

Przypomnijmy uzasadnienie normalnej generacji:
wypełnienie meridionalne zewnętrza $E_K$
odtwarza $S^3$. Włączenie brzegu torusa pełnego
indukuje surjekcję na jego grupę fundamentalną,
z jądrem generowanym przez meridian.
W wolnym iloczynie z amalgamacją można więc
wyeliminować generator grupy fundamentalnej torusa
pełnego i pozostaje dokładnie relacja meridionalna.
Twierdzenie van Kampena dla tego wypełnienia daje
$\pi_1(E_K)/\langle\!\langle\mu_K\rangle\!\rangle
\cong\pi_1(S^3)=1$.
\end{proof}

\begin{twierdzenie}[Wniosek topologiczny; wynik Fintushela--Sterna]
\label{thm:tg-ks-homeo}
Niech $X$ będzie gładką, zamkniętą, spójną,
zorientowaną i po prostu spójną czterorozmaitością.
Niech $T\subset X$ będzie torusem o $T^2=0$,
o otoczeniu typu \emph{cusp} i po prostu spójnym
dopełnieniu. Dla standardowej chirurgii
Fintushela--Sterna $X_K$ jest homeomorficzna z $X$
dla każdego węzła $K\subset S^3$.
\end{twierdzenie}

\begin{proof}[Odwołanie do twierdzenia zewnętrznego]
Udowodniliśmy, że $X_K$ jest po prostu spójna i ma
tę samą homologię, charakterystykę Eulera i sygnaturę.
Do zastosowania klasyfikacji Freedmana nie wystarczają
jednak te dane: trzeba jeszcze wykazać izometrię
całkowitej formy przecięcia oraz zgodność niezmiennika
Kirby'ego--Siebenmanna. Izometrię uzyskaliśmy
w uzupełnieniu dowodu po wniosku
\ref{cor:tg-ks-invariants}, a oba niezmienniki
Kirby'ego--Siebenmanna są zerowe, ponieważ obie
rozmaitości są gładkie. Z twierdzenia Freedmana
wynika homeomorfizm zachowujący orientację.
W tym argumencie nie korzystamy z hipotezy
otoczenia typu \emph{cusp}, którą zachowujemy
dla wspólnego zakresu zastosowań z formułą
Fintushela--Sterna.
\end{proof}

\begin{uwaga}[Różne hipotezy mają różne role]
Prosta spójność dopełnienia pozwala zastosować
van Kampena. Warunek $b_2^+(X)>1$ będzie potrzebny
dopiero do prostej, bezkomorowej formuły Seiberga--Wittena.
Natomiast otoczenie typu \emph{cusp} jest naturalnym
założeniem zapewniającym standardową wersję twierdzenia
Fintushela--Sterna o niezmiennikach gładkich.
Zauważmy jeszcze, że przy $\pi_1(A)=1$ klasa
$[T]$ jest automatycznie prymitywna i niezerowa.
Istotnie, ciąg dokładny pary $(X,A)$ zawiera
\[
H_2(X;\mathbb Z)\longrightarrow H_2(X,A;\mathbb Z)
\longrightarrow H_1(A;\mathbb Z)=0.
\]
Przez wycięcie (z użyciem kołnierzy) i izomorfizm
Thoma dla zorientowanej wiązki normalnej otrzymujemy
$H_2(X,A;\mathbb Z)\cong H_0(T;\mathbb Z)\cong\mathbb Z$.
Pierwsza strzałka jest liczbą przecięcia z $T$:
klasa dysku normalnego jest generatorem względnej
homologii, a każdy punkt przecięcia reprezentuje
taką klasę ze znakiem przecięcia. Surjektywność daje
klasę $x$ z $x\cdot[T]=1$, co wyklucza
zerowość lub podzielność $[T]$ w wolnej homologii.
Przez izometrię wyprowadzoną powyżej to samo
dotyczy $[T_K]$.
\end{uwaga}

\section{Wielomian Alexandera w formule Seiberga--Wittena}
\label{sec:tg-ks-sw}

Niech $X$ będzie zamkniętą, po prostu spójną czterorozmaitością
z $b_2^+(X)>1$. Niezmiennik Seiberga--Wittena można zapisać
jako skończony element pierścienia grupowego
$\mathbb Z[H_2(X;\mathbb Z)/\operatorname{Tor}]$:
\[
\mathcal{SW}_X=\sum_{\beta}SW_X(\beta)\,\exp(\beta),
\]
gdzie suma biegnie po klasach podstawowych, a
$\exp(\beta)$ jest formalnym symbolem o własności
$\exp(\beta)\exp(\gamma)=\exp(\beta+\gamma)$.
Nie jest to funkcja wykładnicza w sensie analizy.
Identyfikację odpowiednich grup homologii $X$
i $X_K$ ustalamy zgodnie z konstrukcją chirurgiczną.
Ponieważ $X$ jest po prostu spójna, dualność
Poincarégo i uniwersalne twierdzenie o współczynnikach
dają wolną grupę $H^2(X;\mathbb Z)$. Klasa wyznacznikowa
wyznacza więc strukturę $\mathrm{Spin}^c$ jednoznacznie:
zmiana struktury o klasę $a\in H^2(X;\mathbb Z)$
zmienia klasę wyznacznikową o $2a$.
Dla wymiaru oczekiwanego $d=0$ współczynnik $SW_X(\beta)$
jest zorientowanym zliczeniem z poprzednich rozdziałów.
W ogólności, przy $b_1=0$ i $d=2r\ge0$, przyjmujemy
całkowitoliczbowy niezmiennik określony przez ewaluację
$r$-tej potęgi klasy punktu (klasy stopnia dwa na
przestrzeni konfiguracji nieredukowalnych modulo
cechowanie) na zorientowanej klasie fundamentalnej
regularnej przestrzeni moduli. Przy $d<0$ lub
nieparzystym $d$ współczynnik przyjmujemy równy zero.
Tę ogólniejszą definicję oraz jej niezmienniczość
i skończoność nośnika przyjmujemy z teorii cechowania
(por. Fintushel--Stern, \emph{Six lectures on four
4-manifolds}, wykład 2, s.~14--15).
Typ prosty oznacza, że wszystkie niezerowe
współczynniki odpowiadają $d=0$; w przykładach K3
i ich modyfikacjach wystarczy ten przypadek.

\begin{twierdzenie}[Fintushel--Stern: formuła \emph{knot surgery}]
\label{thm:tg-ks-formula}
Niech $X$ będzie gładką, zamkniętą, spójną,
zorientowaną i po prostu spójną czterorozmaitością
z $b_2^+(X)>1$. Niech $T\subset X$ będzie homologicznie
nietrywialnym, gładko zanurzonym torusem o $T^2=0$,
będącym regularnym włóknem otoczenia typu
\emph{cusp} (z ramowaniem tego rozwłóknienia) i mającym
po prostu spójne dopełnienie. Niech
$K\subset S^3$ będzie zorientowanym węzłem,
a $\Delta_K(z)$ niech będzie symetrycznie
znormalizowanym wielomianem Alexandera, tj.
$\Delta_K(z^{-1})=\Delta_K(z)$ i $\Delta_K(1)=1$.
Wówczas dla standardowej chirurgii z definicji
\ref{def:tg-ks-operation} zachodzi
\[
\mathcal{SW}_{X_K}=\mathcal{SW}_X\,\Delta_K(u),
\qquad u:=\exp(2[T_K]).
\]
Przez dualność Poincarégo utożsamiamy klasy
wyznacznikowe struktur $\mathrm{Spin}^c$ z klasami
drugiej homologii. Orientację homologiczną
na $X_K$ dobieramy zgodnie z identyfikacją
z $X$; bez tego wyboru równość jest określona
z dokładnością do wspólnego znaku.
Klasa $[T_K]$ odpowiada $[T]$ w izomorfizmie
opisanym po wniosku \ref{cor:tg-ks-invariants}.
Źródło: Fintushel--Stern, \emph{Knots, links, and 4-manifolds},
twierdzenie 1.5, arXiv:dg-ga/9612014v2, s.~4.
\end{twierdzenie}

\begin{uwaga}[Status dowodu i podwójna zmienna]
Twierdzenie Fintushela--Sterna jest \emph{głębokim
wynikiem zewnętrznym}. Jego dowód opiera się na teorii
Seiberga--Wittena, odpowiednich względnych niezmiennikach
i relacjach skein. Nie wynika z ciągu Mayera--Vietorisa.
Niektóre źródła zapisują formułę jako
$\mathcal{SW}_{X_K}=\mathcal{SW}_X\Delta_K(t^2)$,
używając $t=\exp([T])$. W naszym zapisie
$u=\exp(2[T])$, więc są to dokładnie te same równania.
\end{uwaga}

\begin{stwierdzenie}[Klasy podstawowe po chirurgii]
\label{prop:tg-ks-basic}
W hipotezach twierdzenia \ref{thm:tg-ks-formula}
załóżmy dodatkowo, że $\mathcal{SW}_X=1$.
Jeżeli $\Delta_K(z)=\sum_{j=-g}^{g}a_jz^j$,
to $\mathcal{SW}_{X_K}=\sum_{j=-g}^{g}a_j\exp(2j[T_K])$.
W szczególności niezerowe współczynniki $a_j$
określają klasy podstawowe $2j[T_K]$, gdy
$[T_K]$ ma nieskończony rząd; licząc różne wykładniki,
uzyskujemy dokładnie liczbę niezerowych współczynników.
\end{stwierdzenie}

\begin{proof}
Wstawiamy rozwinięcie Laurenta wielomianu Alexandera
do formuły Fintushela--Sterna. Formalne mnożenie
w pierścieniu grupowym jest rozdzielne względem dodawania.
Jeżeli $[T_K]$ ma nieskończony rząd, klasy $2j[T_K]$
dla różnych $j$ są różne, a zatem nie następuje
łączenie odpowiadających im współczynników.
\end{proof}

\begin{uwaga}[Zachowanie wymiaru w formule chirurgii]
Przesunięcie klasy wyznacznikowej $c$ o $2j[T]$
(w identyfikacji Poincarégo) zmienia jej kwadrat o
$4j(c\cdot[T])+4j^2[T]^2=4j(c\cdot[T])$.
Bilans \ref{cor:tg-ks-invariants} oraz wzór
\ref{thm:tg-sw-dimension} dają zatem zmianę wymiaru
oczekiwanego równą $j(c\cdot[T])$ przy identyfikacji
form przecięcia. W szczególnym przypadku
$\mathcal{SW}_X=1$ z poprzedniego stwierdzenia
wszystkie otrzymane klasy $2j[T_K]$ mają kwadrat zero.
Jeżeli ponadto $X$ ma typ prosty, jak w przykładzie K3,
niezerowy niezmiennik klasy zerowej na $X$ wymusza
jej zerowy wymiar oczekiwany, a więc również zerowy
wymiar oczekiwany wszystkich otrzymanych klas. Bez tego założenia
wymiary wszystkich tych klas są po prostu równe
wymiarowi klasy zerowej na $X$.
\end{uwaga}

\section{Przykłady: trójlistnik, węzeł ósemkowy i K3}
\label{sec:tg-ks-examples}

Rozważmy powierzchnię eliptyczną $E(2)$,
czyli powierzchnię K3. Włókno regularne
eliptycznej fibracji jest torusem o samoprzecięciu
zero. Wybieramy standardową fibrację eliptyczną
$E(2)\to\mathbb{CP}^1$ z sekcją i włóknem
ostrzowym, przyjmując istnienie takiej fibracji
z teorii powierzchni eliptycznych. Regularne włókno
$F$ wybieramy w otoczeniu tego włókna osobliwego.
W tym modelu po małym rozszczepieniu osobliwości
ostrzowej na dwa włókna węzłowe dostajemy dwa dyski
zanikające o brzegach reprezentujących bazę
$\pi_1(F)\cong\mathbb Z^2$ (standardowy model
Lefschetza osobliwości ostrzowej).
Oba dyski można obciąć przy brzegu $\nu F$,
otrzymując dyski w dopełnieniu zabijające
dwa kierunki włókna. Trzeci kierunek, meridian,
ogranicza dysk powstały z sekcji $S\cong S^2$
po usunięciu wnętrza dysku $S\cap\nu F$.
W istocie w tych współrzędnych sekcja ustala
punkt włókna; nawet przy zmianie ramowania
ewentualne składowe w kierunkach włókna
już znikają w grupie dopełnienia.
Oznaczmy $A=E(2)\setminus\operatorname{int}\nu F$.
Z van Kampena dla $E(2)=A\cup\nu F$ i surjekcji
$\pi_1(\partial\nu F)\to\pi_1(\nu F)$ wynika,
że $\pi_1(A)$ modulo domknięcie normalne
obrazu meridianu jest $\pi_1(E(2))=1$.
Ponieważ meridian już znika w $\pi_1(A)$,
otrzymujemy $\pi_1(A)=1$. Do samego tego
wniosku wystarcza sekcja sferyczna; dyski
zanikające wyjaśniają dodatkowo lokalny model
używany przy hipotezie typu \emph{cusp}.
Włókno $F$ ma klasę pierwotną w
$H_2(E(2);\mathbb Z)$: przekrój fibracji przecina
włókno z liczbą przecięcia jeden, więc gdyby
$[F]=m\alpha$ z $|m|>1$, otrzymalibyśmy
$1=[F]\cdot[S]=m(\alpha\cdot[S])$, sprzeczność.
W szczególności wykładniki $2j[F]$ dla różnych $j$
są różne w wolnej drugiej homologii.
Dla $E(2)$ przyjmujemy klasyczne dane topologiczne
$\pi_1(E(2))=1$, $\chi(E(2))=24$, $\sigma(E(2))=-16$.
Z dualności Poincarégo mamy $b_2=22$,
a więc $b_2^+=(22-16)/2=3$.
Wynik $K_{E(2)}=0$ i obliczenie
$\mathcal{SW}_{E(2)}=1$ (dla odpowiedniej orientacji
homologicznej) przyjmujemy z teorii powierzchni K3
i nierówności adjunkcji; por. Fintushel--Stern,
\emph{Six lectures on four 4-manifolds},
wykład 2, s.~15 oraz wykład 3, s.~19--20.
Niezerowość klasy zerowej nie wystarczyłaby sama
do wykluczenia pozostałych klas podstawowych.
Stosując pełny wynik i formułę chirurgiczną, dostajemy
\[
\mathcal{SW}_{E(2)_K}=\Delta_K(u),
\qquad u=\exp(2[T_K]).
\]
Wszystkie występujące tu klasy mają kwadrat zero,
gdyż $[T_K]^2=[F]^2=0$. Charakterystyka Eulera
i sygnatura nie zmieniają się, więc dla każdej
klasy $2j[T_K]$ mamy
$d=((2j[T_K])^2-2\cdot24-3(-16))/4=0$.
W szczególności otrzymane rozmaitości mają typ prosty
w sensie Seiberga--Wittena, także dla węzłów
niewłóknistych.

\begin{przyklad}[Chirurgia trójlistnika]
Dla trójlistnika $T_{2,3}$ mamy
$\Delta_{T_{2,3}}(z)=z^{-1}-1+z$. Wobec tego
\[
\mathcal{SW}_{E(2)_{T_{2,3}}}=u^{-1}-1+u.
\]
Rozmaitość $E(2)_{T_{2,3}}$ ma trzy klasy podstawowe
$-2[T_K]$, $0$ i $2[T_K]$ o współczynnikach
odpowiednio $1,-1,1$. Natomiast K3 ma tylko
jedną klasę podstawową, mianowicie zero.
Liczba klas podstawowych jest niezmiennikiem
dyfeomorfizmu zachowującego orientację, dlatego
rozmaitości te nie są dyfeomorficzne w tej kategorii.
Dyfeomorfizm odwracający orientację także nie istnieje:
obie sygnatury wynoszą $-16$, a odwrócenie
orientacji zmieniałoby znak sygnatury.
Ta sama uwaga dotyczy wszystkich dalszych porównań
rozmaitości homeomorficznych z K3. Na mocy twierdzenia
\ref{thm:tg-ks-homeo} są natomiast homeomorficzne.
\end{przyklad}

\begin{przyklad}[Węzeł ósemkowy]
Dla węzła ósemkowego $4_1$ mamy
$\Delta_{4_1}(z)=3-z-z^{-1}$.
Zatem
$\mathcal{SW}_{E(2)_{4_1}}=3-u-u^{-1}$.
Oba omawiane węzły mają rodzaj jeden, lecz
ich wielomiany Alexandera są różne.
W konsekwencji powstałe czterorozmaitości mają
inny multizbiór bezwzględnych wartości niezerowych
współczynników SW. To wystarcza do wykazania,
że nie są dyfeomorficzne, bez konieczności ustalania
konkretnego izomorfizmu ich drugich homologii.
\end{przyklad}

\begin{center}
\begin{tikzpicture}[>=Stealth,line cap=round,line join=round]
 \node[draw=tgBlue,fill=tgBlue!10,rounded corners=4pt,
       text width=3cm,align=center,inner sep=5pt] (k3) at (0,1.2)
       {\small $E(2)$ (K3)\\ \footnotesize $\mathcal{SW}=1$};
 \node[draw=tgTeal,fill=tgTeal!12,rounded corners=4pt,
       text width=3.4cm,align=center,inner sep=5pt] (tr) at (-3.4,-1.1)
       {\small chirurgia z trójlistnikiem\\ \footnotesize $u^{-1}-1+u$};
 \node[draw=tgCoral,fill=tgCoral!10,rounded corners=4pt,
       text width=3.4cm,align=center,inner sep=5pt] (f) at (3.4,-1.1)
       {\small chirurgia z węzłem ósemkowym\\ \footnotesize $3-u-u^{-1}$};
 \draw[-{Stealth[length=2.5mm]},tgPurple,line width=1.7pt] (k3.south west)--(tr.north);
 \draw[-{Stealth[length=2.5mm]},tgPurple,line width=1.7pt] (k3.south east)--(f.north);
 \node[font=\footnotesize,text=tgNavy] at (0,-2.12)
       {Wspólny punkt wyjścia; różne wielomiany SW.};
\end{tikzpicture}
\end{center}

\begin{stwierdzenie}[Suma spójna węzłów a mnożenie wielomianów]
\label{prop:tg-ks-connected-sum}
Niech $K,J\subset S^3$ będą zorientowanymi
węzłami, a $\Delta_K,\Delta_J$ ich
symetrycznie znormalizowanymi
wielomianami Alexandera.
Wówczas
\[
\Delta_{K\#J}(z)=\Delta_K(z)\Delta_J(z).
\]
Jeżeli chirurgię Fintushela--Sterna wykonujemy
na regularnym włóknie eliptycznym K3, to
\[
\mathcal{SW}_{E(2)_{K\#J}}
=\Delta_K(u)\Delta_J(u),
\qquad u=\exp(2[T_{K\#J}]).
\]
\end{stwierdzenie}

\begin{proof}
Wybierzmy powierzchnie Seiferta dla $K$ i $J$
oraz bazy ich pierwszych homologii.
Powierzchnię Seiferta dla sumy spójnej węzłów
otrzymujemy jako brzegową sumę spójną wybranych
powierzchni:
łączymy powierzchnie jednym pasmem przyczepionym
wzdłuż dwóch odcinków brzegu, po jednym na każdej
powierzchni. Każda z powierzchni wraz z odpowiednią
połową pasma deformacyjnie retraktuje się na
pierwotną powierzchnię, a obie pogrubione części
przecinają się w kontraktowalnym prostokącie.
Ciąg Mayera--Vietorisa daje sumę prostą pierwszych
grup homologii, ponieważ homologia stopnia jeden
przecięcia znika, a mapa z jego $H_0$ do sumy
grup $H_0$ obu części jest injektywna. Krzywe bazowe
pochodzące z różnych składników można
wybrać w rozłącznych kulach z wyjątkiem
małego obszaru łączenia, zatem
mieszane liczby splecenia z dodatnimi
przesunięciami wynoszą zero.
Macierz Seiferta otrzymanej powierzchni
ma więc postać blokową $V_K\oplus V_J$.

W konwencji sekcji \ref{sec:kt-alexander}
wielomian Alexandera jest wyznacznikiem
macierzy $zV-V^T$, z dokładnością do
mnożenia przez jednostkę $\pm z^m$.
Równoważny zapis $V-zV^T$ ma ten sam wyznacznik:
jest to macierz $-(zV-V^T)^T$, a rozmiar
macierzy Seiferta węzła jest parzysty.
Wyznacznik macierzy blokowej
jest iloczynem wyznaczników bloków.
Po symetrycznej normalizacji
i warunku $\Delta(1)=1$
otrzymujemy pierwszą równość.
Druga jest zastosowaniem twierdzenia
Fintushela--Sterna
\ref{thm:tg-ks-formula} dla $E(2)$,
którego wielomian Seiberga--Wittena
wynosi $1$.
\end{proof}

\begin{przyklad}[Dwa trójlistniki i węzeł $T(2,5)$]
\label{ex:tg-ks-connected-trefoils}
Niech $T=T(2,3)$ będzie dodatnim trójlistnikiem.
Wtedy
\[
\Delta_{T\#T}(z)
=(z^{-1}-1+z)^2
=z^{-2}-2z^{-1}+3-2z+z^2.
\]
Wzór chirurgiczny na K3 daje więc
pięć klas podstawowych, o
współczynnikach $1,-2,3,-2,1$.
Natomiast węzeł torusowy $T(2,5)$ ma
\[
\Delta_{T(2,5)}(z)
=z^{-2}-z^{-1}+1-z+z^2
\]
i również daje pięć klas podstawowych,
ale wszystkie współczynniki wynoszą $\pm1$.
Te dwie czterorozmaitości są
homeomorficzne z K3 i nie są
dyfeomorficzne, ponieważ
\emph{multizbiory bezwzględnych wartości
niezerowych niezmienników SW} są różne:
odpowiednio $\{1,1,2,2,3\}$
i $\{1,1,1,1,1\}$.
Sama liczba klas podstawowych nie
rozróżniłaby tej pary.
\end{przyklad}

\begin{twierdzenie}[Rodziny o tej samej liczbie klas podstawowych]
\label{thm:tg-ks-equal-basic-count}
Niech $T=T(2,3)$ będzie dodatnim trójlistnikiem.
Dla każdej liczby całkowitej $n\ge2$ zdefiniujmy
dwa węzły $K_n:=\#^{\,n}T$ oraz
$J_n:=T(2,2n+1)$, a przez
$X_n:=E(2)_{K_n}$ i
$Y_n:=E(2)_{J_n}$ oznaczmy wyniki
standardowej chirurgii Fintushela--Sterna
na regularnym włóknie powierzchni K3.
Obie rozmaitości są homeomorficzne z K3
i mają dokładnie $2n+1$ klas podstawowych
Seiberga--Wittena, lecz nie są
dyfeomorficzne między sobą.
\end{twierdzenie}

\begin{proof}
Z multiplikatywności
\ref{prop:tg-ks-connected-sum}
i wielomianu trójlistnika mamy
$\Delta_{K_n}(z)=(z^{-1}-1+z)^n$.
Niech $c_j$ będzie współczynnikiem
przy $z^j$ w tym wielomianie.
Wybierając w każdym z $n$
czynników kolejno składniki
$z^{-1}$, $-1$ lub $z$,
otrzymujemy wykładnik
$j=n-2r-s$, gdzie
$r$ jest liczbą wybranych
czynników $z^{-1}$, a $s$
liczbą wybranych czynników $-1$.
Zatem każdy składnik
współczynnika $c_j$
ma znak $(-1)^s=(-1)^{n-j}$.
Dla $0\le j\le n$ wystarczy wybrać $j$ składników
$z$ i $n-j$ składników $-1$; dla $-n\le j<0$
wybieramy $-j$ składników $z^{-1}$ i $n+j$
składników $-1$. Zatem dla każdego $j$
istnieje składnik o wykładniku $j$ i wszystkie
składniki tego współczynnika mają ten sam znak,
więc $c_j\ne0$.
Są zatem dokładnie
$2n+1$ niezerowe współczynniki.

Współczynnik przy $z^{n-1}$
powstaje, gdy w dokładnie jednym
z $n$ czynników wybierzemy
$-1$, a w pozostałych $z$:
wynosi on $c_{n-1}=-n$.
Z kolei wielomian węzła
$J_n=T(2,2n+1)$ ma postać
$\Delta_{J_n}(z)
=\sum_{j=-n}^{n}(-1)^{n-j}z^j$,
a więc dokładnie $2n+1$
niezerowych współczynników,
każdy równy $\pm1$.

Włókno eliptyczne $F$
jest pierwotną klasą
w $H_2(E(2);\mathbb Z)$.
W standardowej operacji
klasa $[F]$ odpowiada
$[T_K]$, dlatego wykładniki
$2j[T_K]$ są parami różne.
Wzór Fintushela--Sterna daje
dokładnie wskazane liczby
klas podstawowych.

Dla $n\ge2$ współczynnik
$-n$ w wielomianie
$\mathcal{SW}_{X_n}$
ma wartość bezwzględną
większą od jedności,
podczas gdy wszystkie
współczynniki
$\mathcal{SW}_{Y_n}$
mają moduł jeden.
Multizbiór modułów
niezerowych współczynników
niezmiennika Seiberga--Wittena
jest zachowany przy
dyfeomorfizmach, więc
$X_n$ i $Y_n$ nie mogą
być dyfeomorficzne.
Twierdzenie~\ref{thm:tg-ks-homeo}
zapewnia ich homeomorficzność
z K3.
\end{proof}

\section{Nieskończenie wiele egzotycznych struktur na K3}
\label{sec:tg-ks-infinite}

\begin{twierdzenie}[Rodzina egzotycznych struktur]
\label{thm:tg-ks-infinite}
Dla każdego $n\ge1$ niech $K_n=T(2,2n+1)$ będzie
węzłem torusowym o parametrach $(2,2n+1)$,
a $X_n:=E(2)_{K_n}$ wynikiem standardowej chirurgii
Fintushela--Sterna na regularnym włóknie eliptycznym.
Wtedy wszystkie $X_n$ są homeomorficzne z powierzchnią
K3, żadna nie jest dyfeomorficzna z K3 i dla
$n\ne m$ rozmaitości $X_n,X_m$ nie są dyfeomorficzne.
\end{twierdzenie}

\begin{proof}
Wielomian Alexandera węzła torusowego
$T(2,2n+1)$ wyznaczymy ze standardowego wzoru
$\Delta_{T(p,q)}(z)=z^{-(p-1)(q-1)/2}
(1-z^{pq})(1-z)/((1-z^p)(1-z^q))$
dla względnie pierwszych dodatnich $p,q$.
Dla $p=2$, $q=2n+1$ po skróceniu czynników
$1-z$ i $1-z^{2n+1}$ otrzymujemy:
\[
\Delta_{K_n}(z)
 =\frac{z^{2n+1}+1}{z+1}\,z^{-n}
 =\sum_{j=-n}^{n}(-1)^{n-j}z^j.
\]
Pierwsza równość jest tożsamością w pierścieniu
$\mathbb Z[z,z^{-1}]$: dzielenie jest dokładne,
ponieważ $2n+1$ jest nieparzyste. Wzór po prawej
wynika z rozwinięcia
$(z^{2n+1}+1)/(z+1)
=z^{2n}-z^{2n-1}+\cdots-z+1$.
Każdy z $2n+1$ współczynników jest równy $\pm1$
i jest niezerowy.

Twierdzenie Fintushela--Sterna daje
$\mathcal{SW}_{X_n}=\Delta_{K_n}(u)$.
Wynika stąd, że $X_n$ ma dokładnie $2n+1$
klas podstawowych. Ponieważ $2n+1>1$,
nie są dyfeomorficzne ze standardową K3,
która ma jedną klasę podstawową.
Dla $n\ne m$ liczby te również są różne,
więc rozmaitości nie są dyfeomorficzne między sobą.
Jak poprzednio, odwrócenie orientacji jest
wykluczone przez niezerową sygnaturę $-16$.
Twierdzenie \ref{thm:tg-ks-homeo} daje jednocześnie
homeomorfizmy $X_n\cong_{\mathrm{TOP}}E(2)$.
\end{proof}

\begin{uwaga}[Symplektyczność i włóknistość]
Węzły torusowe $T(2,2n+1)$ są włókniste
(klasyczny wynik o węzłach torusowych).
Jeśli $K$ jest włóknisty, jego zewnętrze jest
wiązką nad $S^1$ o włóknie będącym powierzchnią
Seiferta z jednym brzegiem $\lambda_K$.
Chirurgia Dehna o współczynniku zero polega na
dołączeniu torusa pełnego wzdłuż tej długości.
Domykamy w ten sposób każde włókno dyskiem,
a monodromię, ustaloną przy brzegu, przedłużamy
przez identyczność na dysku. Powstała trójrozmaitość
$M_K$ jest wiązką domkniętych powierzchni nad $S^1$.
Rdzeń dołączonego torusa pełnego, oznaczony $m$,
jest sekcją tej wiązki, a
$M_K\setminus\operatorname{int}\nu m\cong E_K$.
W $S^1\times M_K$ torus $S^1\times m$ jest
sekcją wiązki powierzchni nad $T^2$ i ma
trywialną wiązkę normalną w tym modelu.
Konstrukcja Thurstona daje formę symplektyczną
na tej wiązce (klasa kohomologii dodatnia
na włóknach istnieje, ponieważ sekcja
$m$ przecina włókno raz). Dodając dostatecznie
dużą wielokrotność formy pola z bazy $T^2$,
można zarazem uczynić torus sekcji symplektycznym.
Po uzgodnieniu pól torusów przez przeskalowanie
form stosujemy twierdzenie Gompfa o symplektycznej
sumie włóknistej wzdłuż $F$ i $S^1\times m$.
Usunięcie otoczenia drugiego torusa zostawia
$S^1\times E_K$, a jego meridian odpowiada
$\lambda_K$ (brzegowi dysku domykającego
powierzchnię Seiferta). Z ramowaniem produktowym
otrzymujemy właśnie standardowe sklejenie
\ref{def:tg-ks-operation}. W konsekwencji
wszystkie $X_n$ mają struktury symplektyczne.
Twierdzenia Thurstona i Gompfa przyjmujemy
jako wyniki zewnętrzne. W szczególności jest to
rodzina \emph{symplektycznych} struktur egzotycznych
na topologicznej powierzchni K3.
\end{uwaga}

\section{Czego wielomian Alexandera nie wykrywa?}
\label{sec:tg-ks-limits}

\begin{stwierdzenie}[Ograniczenie formuły]
\label{prop:tg-ks-limit}
Niech $K_1,K_2\subset S^3$ będą dwoma węzłami,
dla których symetrycznie znormalizowane wielomiany
Alexandera są równe:
$\Delta_{K_1}=\Delta_{K_2}$.
Jeśli $X$ i $T$ spełniają hipotezy twierdzenia
\ref{thm:tg-ks-formula}, to po odpowiednich
identyfikacjach homologii zachodzi
$\mathcal{SW}_{X_{K_1}}=\mathcal{SW}_{X_{K_2}}$.
Nie wynika stąd, że $X_{K_1}$ i $X_{K_2}$
są dyfeomorficzne.
\end{stwierdzenie}

\begin{proof}
Pierwsza teza jest natychmiastowym podstawieniem
do wzoru Fintushela--Sterna. Druga wypowiedź oznacza
jedynie ograniczenie logiczne metody: równość tych
niezmienników nie jest twierdzeniem klasyfikacyjnym
wymuszającym dyfeomorfizm.
\end{proof}

\begin{przyklad}[Wielomian trywialny nie oznacza węzła trywialnego]
Istnieją nietrywialne węzły $K$ z $\Delta_K(z)=1$,
na przykład węzeł Kinoshity--Terasaki (Kinoshita--Terasaka, $11n42$).
Dla takich węzłów w formule z K3 otrzymujemy
$\mathcal{SW}_{E(2)_K}=1=\mathcal{SW}_{E(2)}$.
Jest to równość obliczonych niezmienników, a nie
dowód dyfeomorficzności. Z drugiej strony pokazuje,
że sam wielomian Alexandera nie może rozstrzygnąć
wszystkich pytań o strukturę gładką.
\end{przyklad}

\section{Porównanie z wymierną kontrakcją Parka}
\label{sec:tg-ks-compare}

Operacje wymiernej kontrakcji i chirurgii wzdłuż węzła
prowadzą do egzotycznych struktur gładkich przy
odmiennych założeniach. Dla Parka punktem wyjścia
jest ujemnie określony łańcuch sfer o samoprzecięciach
różnych od zera, zastępowany wymierną kulą homologii.
W konstrukcji Fintushela--Sterna modyfikujemy otoczenie
torusa o samoprzecięciu zero.

\begin{center}
\begin{tikzpicture}[>=Stealth,line cap=round,line join=round]
  \node[draw=tgCoral,fill=tgCoral!11,rounded corners=4pt,
        text width=4.2cm,align=center,inner sep=7pt] (a) at (-3.1,0)
        {\small wymierna kontrakcja\\
         \footnotesize łańcuch sfer $C_p$\\
         \footnotesize $b_2^+$ może wynosić $1$};
  \node[draw=tgTeal,fill=tgTeal!11,rounded corners=4pt,
        text width=4.2cm,align=center,inner sep=7pt] (b) at (3.1,0)
        {\small chirurgia wzdłuż węzła\\
         \footnotesize torus $[T]^2=0$\\
         \footnotesize formuła SW: $b_2^+>1$};
  \draw[<->,tgPurple,line width=1.7pt] (a)--(b);
  \node[font=\footnotesize,text=tgNavy] at (0,-1.45)
        {Różne operacje, wspólny cel: wykrywanie typu DIFF.};
\end{tikzpicture}
\end{center}

Nie wolno bezpośrednio zastosować twierdzenia
\ref{thm:tg-ks-formula} do rozmaitości Parka $X_7$,
ponieważ $b_2^+(X_7)=1$. W tym przypadku
niezbędne jest rozróżnienie komór i osobne
argumenty analityczne omówione wcześniej.
Ponadto sama obecność zanurzonego torusa o
$T^2=0$ nie gwarantuje istnienia otoczenia typu
\emph{cusp} ani prostej spójności dopełnienia.

\begin{uwaga}[Co jest dowiedzione, a co przyjęte]
W tym rozdziale samodzielnie przeprowadziliśmy
obliczenia homologii zewnętrza węzła i części
$S^1\times E_K$, dowód prostej spójności
$X_K$ przy trywialnej grupie dopełnienia torusa
oraz rachunki wielomianowe dla wielu przykładów.
Korzystamy zewnętrznie z addytywności sygnatur,
twierdzenia Freedmana, formuły Fintushela--Sterna
i konstrukcji symplektycznej sumy włóknistej.
Izometrię form przecięcia wyprowadziliśmy
z odwzorowania stopnia jeden. To formuła
Fintushela--Sterna umożliwia rozróżnianie
struktur gładkich, a suma włóknista zapewnia
symplektyczność wskazanej rodziny.
\end{uwaga}


% ITERACJA 16: synteza rachunków grafowych, chirurgicznych i cechowania
